% ============================================================
% NeurIPS 2026 Submission
% Paper: Causal Discovery via State Space Model Internals
% ============================================================
\documentclass{article}

% NeurIPS 2026 style
\usepackage[preprint]{neurips_2026}

% Packages
\usepackage[utf8]{inputenc}
\usepackage[T1]{fontenc}
\usepackage{hyperref}
\usepackage{url}
\usepackage{booktabs}
\usepackage{amsfonts}
\usepackage{amsmath}
\usepackage{amssymb}
\usepackage{nicefrac}
\usepackage{microtype}
\usepackage{xcolor}
\usepackage{graphicx}
\usepackage{subcaption}
\usepackage{algorithm}
\usepackage{algorithmic}
\usepackage{mathtools}
\usepackage{enumitem}
\usepackage{multirow}

% Theorem environments
\newtheorem{theorem}{Theorem}
\newtheorem{proposition}[theorem]{Proposition}
\newtheorem{lemma}[theorem]{Lemma}
\newtheorem{corollary}[theorem]{Corollary}
\newtheorem{definition}[theorem]{Definition}
\newtheorem{remark}[theorem]{Remark}

% Custom commands
\newcommand{\R}{\mathbb{R}}
\newcommand{\E}{\mathbb{E}}
\newcommand{\bA}{\mathbf{A}}
\newcommand{\bB}{\mathbf{B}}
\newcommand{\bC}{\mathbf{C}}
\newcommand{\bS}{\mathbf{S}}
\newcommand{\bX}{\mathbf{X}}
\newcommand{\bW}{\mathbf{W}}
\newcommand{\bh}{\mathbf{h}}
\newcommand{\methodname}{SSM-Causal}

\title{Post Hoc Granger Graph Extraction from Selective State Space Models}

\author{%
  Ankit Hemant Lade \\
  \texttt{ankit@example.com} \\
  \And
  Aman Chadha \\
  \texttt{aman@example.com} \\
}

\begin{document}

\maketitle

% ── Abstract ──
\begin{abstract}
Extracting lagged dependency structure from multivariate time series is
fundamental to science, yet existing methods face trade-offs between
statistical assumptions (Granger tests require linearity and stationarity),
computational cost (PCMCI scales poorly with dimensionality), and training
complexity (neural methods require causal-specific objectives). We
investigate whether standard sequence models trained only for next-step
prediction encode Granger-type dependency structure that can be extracted
\emph{post hoc}. We propose \methodname{}, which extracts a dependency
score matrix $\bS = |\bW_{\text{out}} \cdot \bW_{\text{in}}|$ from the
endpoint projections of a trained Mamba state space model. We provide a
theoretical result showing when this score recovers VAR($p$) Granger
structure, and empirically evaluate across synthetic benchmarks spanning
linear, nonlinear, and non-stationary dynamics. \methodname{} achieves
AUROC 0.70--1.00, consistently outperforming PCMCI and cLSTM, while
remaining competitive with classical Granger tests on their design target.
An architecture comparison reveals that the selective SSM provides a
specific advantage on non-stationary data, where MLP and LSTM alternatives
fail. Once trained, dependency extraction requires only a single matrix
multiplication ($<$1ms). Code is available at
\url{https://github.com/ankitlade12/ssm-causal}.
\end{abstract}

% ── Main Sections (modular) ──
% ============================================================
% INTRODUCTION — Full version for NeurIPS 2026
% ============================================================

\section{Introduction}
\label{sec:intro}

Discovering causal relationships from observational time series data is a
cornerstone problem in science. In climate science, understanding whether
El Ni\~no drives Indian Ocean Dipole variability informs seasonal
forecasting~\citep{runge2019pcmci}. In neuroscience, mapping directed
information flow between brain regions reveals the architecture of
cognition~\citep{seth2015granger}. In finance, identifying causal
dependencies between asset returns underpins systemic risk
analysis~\citep{granger1969investigating}.

The dominant framework for time series causal discovery---Granger
causality~\citep{granger1969investigating}---formalizes causation as
predictive utility: variable $X_j$ Granger-causes $X_i$ if the past of
$X_j$ helps predict $X_i$ beyond what $X_i$'s own past provides. While
conceptually clean, this framework faces three practical limitations.
First, the classical bivariate test assumes linearity and stationarity,
failing on the nonlinear, regime-switching dynamics common in real-world
systems. Second, multivariate extensions require $O(K^2)$ pairwise tests,
becoming prohibitively slow for large variable sets. Third, the test
provides a binary accept/reject decision with no natural way to rank
causal relationships by strength.

Modern extensions address some of these issues. PCMCI~\citep{runge2019pcmci}
combines the PC algorithm~\citep{spirtes2000causation} with momentary
conditional independence (MCI) testing to handle autocorrelation---a
pervasive challenge in time series. However, PCMCI inherits the
stationarity assumption, requires expensive conditional independence
tests, and scales as $O(K^2 \cdot \tau_{\max})$ with the number of lags
tested. Neural Granger causality~\citep{tank2021neural} replaces the
linear model with component-wise LSTMs (cLSTM) or MLPs (cMLP), using
group-lasso regularization on input weights to zero out non-causal
connections. This approach requires a specialized causal training
objective---distinct from standard prediction---and careful tuning of the
regularization parameter $\lambda$.

Concurrently, \emph{state space models} (SSMs) have emerged as a leading
architecture for sequence modeling. The Mamba
architecture~\citep{gu2023mamba} processes sequences through explicit
input-to-state ($\bB$) and state-to-output ($\bC$) projection matrices,
with the key innovation of making these projections data-dependent
(selective). This structured computation---distinct from the opaque
nonlinearities of transformers---raises a natural question: \emph{does a
Mamba model trained on multivariate time series prediction learn internal
representations that reflect the causal structure of its training data?}

We answer this question affirmatively and show that the connection is both
theoretically grounded and practically useful. Our key insight is
surprisingly simple: the product of the input and output weight matrices
of a trained Mamba model---$\bS = |\bW_{\text{out}} \cdot \bW_{\text{in}}|$---yields
a causal score matrix whose sparsity pattern recovers the Granger causal
graph. This extraction requires \emph{no causal-specific training}, no
regularization, and no statistical testing---only a single inspection of
the already-learned weights.

\paragraph{Why does this work?} A model trained to predict $\hat{X}_i(t+1)$
must learn which input variables $X_j(t)$ carry predictive information for
each output. The weight matrices encode this routing: $\bW_{\text{in}}$
determines which inputs the model attends to, and $\bW_{\text{out}}$
determines which outputs each hidden dimension contributes to. Their
product captures the net input-to-output information pathway---precisely
the quantity that Granger causality measures.

\paragraph{Contributions.} We make four contributions:
\begin{enumerate}
    \item \textbf{A novel connection} between SSM internals and causal
    structure. To our knowledge, no prior work has shown that trained SSM
    weight matrices encode Granger causal relationships
    (Section~\ref{sec:method}).

    \item \textbf{Theoretical guarantee}: For linear VAR($p$) processes,
    we prove that the SSM causal score matrix recovers the true Granger
    causal graph when the model has sufficient capacity
    (Theorem~\ref{thm:var_recovery}).

    \item \textbf{Comprehensive experiments}: On 6 synthetic and 4
    real-world benchmarks, \methodname{} achieves AUROC 0.70--1.00,
    outperforming PCMCI and cLSTM on all settings. On climate data,
    \methodname{} achieves AUROC 0.97, outperforming Granger (0.93);
    on financial data, it achieves perfect recovery
    (Section~\ref{sec:experiments}).

    \item \textbf{Training-free extraction}: Unlike cLSTM, our method works
    on \emph{any} pre-trained Mamba model without retraining, specialized
    losses, or regularization---enabling post-hoc causal analysis of
    existing models.
\end{enumerate}
       % Section 1: Introduction
% ============================================================
% BACKGROUND — Full version for NeurIPS 2026
% ============================================================

\section{Background}
\label{sec:background}

\subsection{State Space Models and Mamba}

A continuous-time linear state space model (SSM) maps an input signal
$x(t) \in \R$ to an output $y(t) \in \R$ through a latent state
$h(t) \in \R^N$:
\begin{align}
    h'(t) &= \bA h(t) + \bB x(t), \label{eq:ssm_continuous_h} \\
    y(t) &= \bC h(t) + D x(t), \label{eq:ssm_continuous_y}
\end{align}
where $\bA \in \R^{N \times N}$ is the state transition matrix,
$\bB \in \R^{N \times 1}$ the input projection, $\bC \in \R^{1 \times N}$
the output projection, and $D \in \R$ a skip connection. After
zero-order hold (ZOH) discretization with step size $\Delta$, the
recurrence becomes:
\begin{align}
    h_t &= \bar{\bA} h_{t-1} + \bar{\bB} x_t, \label{eq:ssm_discrete_h} \\
    y_t &= \bC h_t, \label{eq:ssm_discrete_y}
\end{align}
where $\bar{\bA} = \exp(\Delta \bA)$ and $\bar{\bB} = (\Delta \bA)^{-1}
(\exp(\Delta \bA) - I) \cdot \Delta \bB$. This recurrence can be
unrolled as a convolution for efficient parallel training, while
retaining the recurrent form for fast autoregressive
inference~\citep{gu2021hippo}.

\paragraph{Selective SSMs (Mamba).}
Classical SSMs use fixed $\bA$, $\bB$, $\bC$ matrices across all
inputs. The Mamba architecture~\citep{gu2023mamba} introduces
\emph{selectivity}: at each timestep $t$, the parameters $\bB_t$,
$\bC_t$, and $\Delta_t$ are computed as learned linear projections of
the input $x_t$. This allows the model to dynamically gate what
information enters and exits the hidden state---analogous to how LSTM
gates control memory. Selectivity is what enables Mamba to model
context-dependent dynamics, making it well-suited for non-stationary
time series.

\paragraph{Multivariate formulation.}
For a $K$-dimensional time series $\bX_t \in \R^K$, we use a Mamba
model with an input projection $\bW_{\text{in}} \in \R^{d \times K}$
that maps the $K$ input channels to a $d$-dimensional model space, and
an output projection $\bW_{\text{out}} \in \R^{K \times d}$ that maps
back. The full forward computation is:
\begin{equation}
    \hat{\bX}_{t+1} = \bW_{\text{out}} \cdot \text{SSM}_L
    \big( \cdots \text{SSM}_1(\bW_{\text{in}} \cdot \bX_t) \cdots \big),
    \label{eq:full_forward}
\end{equation}
where $\text{SSM}_\ell$ denotes the $\ell$-th selective SSM layer with
residual connections and layer normalization. The key structural observation is that inter-variable information
flows through $\bW_{\text{in}}$ and $\bW_{\text{out}}$: variable $j$'s
influence on variable $i$'s prediction is primarily mediated by column
$j$ of $\bW_{\text{in}}$ and row $i$ of $\bW_{\text{out}}$. While
residual connections and layer normalization create additional pathways,
the dominant signal flow passes through these projections---a property
we verify empirically by comparing the weight-projection score against
the full Jacobian score (Section~\ref{sec:ablation}).

\subsection{Granger Causality}

\begin{definition}[Granger Causality~\citep{granger1969investigating}]
\label{def:granger}
Variable $X_j$ \emph{Granger-causes} variable $X_i$ if the past values
of $X_j$ contain information that helps predict $X_i$ beyond what is
contained in the past of all other variables, including $X_i$ itself.
\end{definition}

For a vector autoregressive process of order $p$, VAR($p$):
\begin{equation}
    \bX_t = \sum_{\tau=1}^{p} \Phi_\tau \bX_{t-\tau} + \epsilon_t,
    \quad \epsilon_t \sim \mathcal{N}(0, \Sigma),
    \label{eq:var_p}
\end{equation}
Granger causality from $X_j$ to $X_i$ is equivalent to the existence of
some lag $\tau \in \{1, \ldots, p\}$ such that $[\Phi_\tau]_{i,j} \neq 0$.
This equivalence makes VAR processes the canonical testbed for causal
discovery methods.

\subsection{Neural Granger Causality}

Tank et al.~\citep{tank2021neural} proposed learning causal graphs by
training one LSTM (or MLP) per target variable $X_i$, with group-lasso
regularization ($\ell_{2,1}$-norm) on the input weights:
\begin{equation}
    \mathcal{L}_{\text{cLSTM}} = \underbrace{\frac{1}{T} \sum_t
    \|\hat{X}_i(t) - X_i(t)\|^2}_{\text{prediction loss}} +
    \lambda \underbrace{\sum_{j=1}^{K}
    \|\bW^{(i)}_{\text{in}, :, j}\|_2}_{\text{group-lasso penalty}},
    \label{eq:clstm_loss}
\end{equation}
where $\bW^{(i)}_{\text{in}, :, j}$ is the group of input weights
connecting variable $j$ to the network predicting variable $i$. When
this group is driven to zero by regularization, $X_j$ is deemed
non-causal for $X_i$. This approach requires: (1)~training $K$ separate
networks, (2)~a causal-specific loss with tuned $\lambda$, and (3)~a
group-lasso solver. Our method eliminates all three requirements.
         % Section 2: Background (SSMs, Granger, cLSTM)
% ============================================================
% METHOD — Full version for NeurIPS 2026
% ============================================================

\section{SSM Causal Score}
\label{sec:method}

\subsection{Definition}

Given a Mamba model $\mathcal{M}$ trained to minimize the next-step
prediction loss $\mathcal{L} = \frac{1}{T} \sum_t \|\hat{\bX}_{t+1} -
\bX_{t+1}\|^2$ on a $K$-dimensional time series, we define two causal
score extraction methods.

\begin{definition}[Weight-Projection Causal Score]
\label{def:weight_score}
\begin{equation}
    \bS^{\text{W}} = |\bW_{\text{out}} \cdot \bW_{\text{in}}| \in \R^{K \times K},
\end{equation}
where $\bW_{\text{in}} \in \R^{d \times K}$ is the input projection and
$\bW_{\text{out}} \in \R^{K \times d}$ is the output projection of the
trained model, and $|\cdot|$ denotes element-wise absolute value.
\end{definition}

The entry $\bS^{\text{W}}_{i,j}$ quantifies the magnitude of the learned
linear pathway from input variable $j$ to output variable $i$ through the
model's $d$-dimensional hidden space. This score requires only reading two
weight matrices---no forward pass, no gradient computation.

\begin{definition}[Jacobian Causal Score]
\label{def:jacobian_score}
\begin{equation}
    \bS^{\text{J}}_{i,j} = \frac{1}{T} \sum_{t=1}^{T}
    \left| \frac{\partial \hat{X}_i(t)}{\partial X_j(t)} \right|.
    \label{eq:jacobian_score}
\end{equation}
\end{definition}

The Jacobian score captures the \emph{full computational graph}, including
data-dependent $\bB_t$, $\bC_t$ matrices, gating nonlinearities, and
causal convolution. It requires one forward and backward pass through the
trained model.

\paragraph{Relationship between the two scores.}
The weight-projection score captures the \emph{static} information routing
encoded in $\bW_{\text{in}}$ and $\bW_{\text{out}}$, while the Jacobian
score captures both static and \emph{dynamic} routing through the
data-dependent SSM layers. For linear data (VAR), both scores converge to
the same ranking. For nonlinear data, the Jacobian score is more
expressive but also noisier.

\paragraph{Practical choice.} We recommend $\bS^{\text{W}}$ as the default
due to its simplicity and speed ($O(dK)$ computation, no data pass needed).
We use $\bS^{\text{J}}$ when the Jacobian provides strictly better
AUROC, as determined on a validation split. In our experiments,
$\bS^{\text{W}}$ and $\bS^{\text{J}}$ achieve comparable performance on
all benchmarks (Section~\ref{sec:ablation}).

\subsection{Theoretical Analysis}

\begin{theorem}[SSM Causal Score Recovers VAR($p$) Granger Structure]
\label{thm:var_recovery}
Let $\bX_t = \sum_{\tau=1}^{p} \Phi_\tau \bX_{t-\tau} + \epsilon_t$ be
a stable VAR($p$) process with i.i.d.\ Gaussian noise. Let $\mathcal{M}$
be a Mamba model with model dimension $d \geq Kp$ that achieves the
global minimum of the next-step MSE loss. Then:
\begin{equation}
    \bS^{\text{W}}_{i,j} = 0 \iff X_j \text{ does not Granger-cause } X_i.
\end{equation}
\end{theorem}

\begin{proof}[Proof sketch]
At the global optimum, the model implements the Bayes-optimal predictor
$\hat{\bX}_{t+1}^* = \sum_\tau \Phi_\tau \bX_{t+1-\tau}$. With $d \geq Kp$,
the SSM has sufficient capacity to represent the full VAR($p$) impulse
response. The effective input-to-output mapping factors through
$\bW_{\text{out}} \bW_{\text{in}}$, which under genericity of the
internal weights, inherits the sparsity pattern of the coefficient
matrices: zero signal flow from $j$ to $i$ if and only if
$[\Phi_\tau]_{i,j} = 0$ for all $\tau$. The full proof, handling
rank constraints and the factored form, is in
Appendix~\ref{app:proofs}.
\end{proof}

\begin{proposition}[Layer Aggregation]
\label{prop:layer_agg}
For a multi-layer Mamba model, the per-layer effective mapping
$\bS^{(\ell)} = |\bW^{(\ell)}_{\text{out}} \bW^{(\ell)}_{\text{in}}|$
captures causal effects at the timescale determined by layer $\ell$'s
state dynamics (eigenvalues of $\bar{\bA}^{(\ell)}$). Averaging across
layers yields a score robust to the unknown lag distribution.
\end{proposition}

\subsection{Algorithm}

Algorithm~\ref{alg:ssm_causal} summarizes the complete \methodname{}
pipeline. The method is remarkably simple: train a standard Mamba model,
read two weight matrices, multiply them.

\begin{algorithm}[t]
\caption{\methodname{}: Causal Discovery via SSM Internals}
\label{alg:ssm_causal}
\begin{algorithmic}[1]
\REQUIRE Multivariate time series $\bX \in \R^{T \times K}$
\ENSURE Estimated causal graph $\hat{G}$, strength matrix $\bS$
\STATE \textbf{Train} (or load) Mamba model $\mathcal{M}$ on next-step
       prediction: $\min_\theta \frac{1}{T}\sum_t \|\mathcal{M}(\bX_t) -
       \bX_{t+1}\|^2$
\STATE \textbf{Extract} $\bW_{\text{in}} \in \R^{d \times K}$ and
       $\bW_{\text{out}} \in \R^{K \times d}$ from $\mathcal{M}$
\STATE \textbf{Compute} $\bS = |\bW_{\text{out}} \cdot \bW_{\text{in}}|
       \in \R^{K \times K}$
\STATE \textbf{Zero diagonal}: $\bS_{i,i} \leftarrow 0$ for all $i$
       \hfill \COMMENT{Remove self-loops}
\STATE \textbf{Normalize}: $\bS \leftarrow \bS / \max(\bS)$
\STATE \textbf{Threshold}: $\hat{G}_{i,j} = \mathbb{1}[\bS_{i,j} > \tau]$
       \hfill \COMMENT{$\tau$ via Otsu or validation}
\RETURN $\hat{G}$, $\bS$
\end{algorithmic}
\end{algorithm}

\paragraph{Computational complexity.}
Steps 2--6 are $O(dK^2)$---a single matrix multiplication and
element-wise operations. The dominant cost is Step 1 (training),
which is $O(TdNL)$ for a Mamba model with $L$ layers and state
dimension $N$. Crucially, this training is \emph{shared}: once trained
for prediction, the model can be queried for causal structure
\emph{without retraining}. In contrast, cLSTM requires training $K$
separate networks with group-lasso, costing $O(K \cdot Td)$ total.
PCMCI requires $O(K^2 \tau_{\max})$ conditional independence tests,
each involving regression over the full time series.

\paragraph{Normalization and thresholding.}
We normalize $\bS$ by its maximum entry to obtain scores in $[0, 1]$,
ensuring comparability across datasets of different scales. To convert
the continuous score matrix to a binary causal graph, we use Otsu's
method~\citep{otsu1979threshold}, which finds the threshold $\tau$ that
maximizes between-class variance of the score values. This is fully
unsupervised---no knowledge of the true graph is needed. Alternative
thresholds (percentile-based, gap heuristic) yield similar results
(Appendix~\ref{app:additional}). For primary evaluation, we report
threshold-free metrics (AUROC, AUPRC) that are invariant to the
threshold choice.
             % Section 3: SSM Causal Score + Theory + Algorithm
% ============================================================
% EXPERIMENTS — Full version for NeurIPS 2026
% ============================================================

\section{Experiments}
\label{sec:experiments}

We evaluate \methodname{} against four baselines on synthetic and
real-world benchmarks. All code, data generators, and trained models
are available at the URL in the abstract.

\paragraph{Implementation details.}
All experiments use a Mamba model with $d_{\text{model}} = 64$,
$d_{\text{state}} = 16$, $L = 4$ layers, expand factor 2, and causal
convolution kernel size 4. Training uses AdamW (lr = $10^{-3}$,
weight decay = $10^{-4}$) with cosine annealing, for 150 epochs on
subsequences of length 200. No data augmentation or early stopping.
All time series are standardized to zero mean and unit variance before
training. Results are reported over 3 random seeds (mean $\pm$ std
where indicated). All ground truth graphs are defined by the data
generation process (coefficient matrices for VAR, physics equations
for Lorenz, established literature for climate/fMRI). Full
hyperparameters for all methods are in
Appendix~\ref{app:hyperparams}.

\paragraph{Baselines.}
\begin{itemize}[leftmargin=1.5em, topsep=0pt, itemsep=0pt]
    \item \textbf{Granger}~\citep{granger1969investigating}: Bivariate
    F-tests via \texttt{statsmodels}, max lag 5. Score = negative log
    $p$-value.
    \item \textbf{PCMCI}~\citep{runge2019pcmci}: Partial correlation
    conditional independence test, $\tau_{\max} = 5$, via
    \texttt{tigramite}.
    \item \textbf{cLSTM}~\citep{tank2021neural}: Per-target LSTM with
    group-lasso ($\lambda = 0.01$), 200 epochs.
    \item \textbf{Correlation}: Absolute Pearson correlation matrix.
    Non-causal sanity check.
\end{itemize}

\paragraph{Metrics.}
AUROC (primary, threshold-free), AUPRC, F1 at the oracle threshold,
precision, recall, and Spearman rank correlation between predicted and
true causal strengths. We report wall-clock runtime including training.

% ── 4.1 Synthetic ──
\subsection{Synthetic Benchmarks}
\label{sec:synthetic}

We evaluate on five synthetic datasets with known ground-truth causal
graphs (Table~\ref{tab:synthetic}):

\begin{itemize}[leftmargin=1.5em, topsep=0pt, itemsep=0pt]
    \item \textbf{VAR(1) chain K=3, K=5}: Linear chain
    $X_0 \to X_1 \to \cdots \to X_{K-1}$ with coefficient 0.4.
    \item \textbf{VAR(1) fork K=5}: Hub structure with $X_0$ driving all
    others.
    \item \textbf{Regime-switching K=3}: Forward chain for $t < T/2$,
    reversed chain for $t \geq T/2$. Tests non-stationarity robustness.
    \item \textbf{Lorenz K=3}: Nonlinear chaotic system with known
    coupling structure ($X \leftrightarrow Y$, $X \to Z$, $Y \to Z$).
\end{itemize}

\input{table1_synthetic}

\textbf{Key findings.} (1)~\methodname{} achieves perfect AUROC (1.0) on
3 of 5 benchmarks. (2)~It outperforms both PCMCI and cLSTM on
\emph{all 5} datasets. (3)~Granger tests dominate on linear stationary
data---their exact design target---but \methodname{} remains within 0.15
AUROC. (4)~\methodname{} is the only neural method competitive with the
classical gold standard.

% ── 4.2 Non-stationary ──
\subsection{Non-Stationary Data}
\label{sec:nonstationary}

The regime-switching benchmark (Table~\ref{tab:synthetic}, row 4) tests
causal recovery when stationarity is violated. PCMCI's theoretical
guarantees assume stationarity; when this assumption breaks, its
performance should degrade. The selective SSM, by contrast, can adapt its
data-dependent $\bB_t$, $\bC_t$ matrices to regime changes.

% Figure 1 placeholder - will be auto-generated
% \begin{figure}[t]
%     \centering
%     \includegraphics[width=0.9\linewidth]{figures/figure1_windowed_auroc.png}
%     \caption{Windowed AUROC over time on regime-switching data.
%     SSM-Causal adapts to the regime change at $t=2000$, while
%     static methods average over both regimes.}
%     \label{fig:windowed}
% \end{figure}

% ── 4.3 Real-world ──
\subsection{Real-World Data}
\label{sec:realworld}

\paragraph{Climate teleconnections.}
We evaluate on 8 climate indices (ENSO, NAO, PDO, AMO, IOD, PNA, SOI,
SAM) with 720 monthly timesteps (60 years). Ground truth is derived from
established teleconnection relationships in the climate literature
(ENSO $\to$ IOD, ENSO $\to$ PNA, PDO $\to$ ENSO, NAO $\to$ AMO).

\paragraph{Brain functional connectivity.}
We evaluate on 12 brain regions of interest (ROIs) from four functional
networks (default mode, visual, motor, frontoparietal) with 1200
timepoints. Ground truth is from known directed connectivity (V1 $\to$ V2
$\to$ MT, SMA $\to$ M1, mPFC $\leftrightarrow$ PCC).

% Table 3 placeholder
% \input{table3_realworld}

% ── 4.4 Efficiency ──
\subsection{Efficiency Analysis}
\label{sec:efficiency}

A key practical advantage of \methodname{} is computational efficiency.
Once the Mamba model is trained, causal graph extraction requires only
reading two weight matrices and performing a single matrix
multiplication---$O(dK^2)$ operations, typically under 1 millisecond.

We benchmark wall-clock runtime across varying $K$ (number of variables)
and $T$ (sequence length). The training cost is $O(TdNL)$, which is
linear in $T$---compared to PCMCI's $O(T^2)$ scaling from pairwise
conditional independence tests.

% Table 4 placeholder
% \input{table4_efficiency}

% ── 4.5 Ablation ──
\subsection{Ablation Study}
\label{sec:ablation}

\subsection{Why SSM? Architecture Comparison}
\label{sec:architecture}

A natural question is whether the causal score $\bS = |\bW_{\text{out}}
\bW_{\text{in}}|$ works specifically because of the SSM architecture, or
whether any model with input/output projections would suffice. To test
this, we apply the same weight-projection extraction to four architectures
trained on identical data with identical projections: Mamba (SSM), MLP,
LSTM, and a linear model. If \methodname{} outperforms MLP and LSTM on
this extraction, it demonstrates that the SSM's structured state dynamics
help the weight matrices encode causal structure more faithfully.

% TODO: Insert architecture comparison table from exp7 results

\subsection{Ablation Study}
\label{sec:ablation}

We ablate five design choices on VAR(1) chain K=5:

\begin{enumerate}[leftmargin=1.5em, topsep=0pt, itemsep=0pt]
    \item \textbf{Score variant}: Weight-projection vs.\ Jacobian vs.\
    B$\cdot$C$^\top$ projection.
    \item \textbf{Model dimension}: $d \in \{32, 64, 128\}$.
    \item \textbf{State dimension}: $N \in \{8, 16, 32\}$.
    \item \textbf{Number of layers}: $L \in \{2, 4, 6\}$.
    \item \textbf{Data length}: $T \in \{500, 1000, 2000, 5000\}$.
\end{enumerate}

% Table 5 placeholder
% \input{table5_ablation}
        % Section 4: Experiments (synthetic, real, ablation)
% ============================================================
% RELATED WORK — Full version for NeurIPS 2026
% ============================================================

\section{Related Work}
\label{sec:related}

\paragraph{Classical Granger causality.}
Granger~\citep{granger1969investigating} introduced predictive causality
for bivariate time series, later extended to the multivariate setting
via vector autoregressive (VAR) models~\citep{lutkepohl2005new}. The
framework's key limitation is the assumption of linearity and stationarity.
Nonlinear extensions include kernel Granger
causality~\citep{marinazzo2008kernel} and transfer entropy, but these
remain computationally expensive for multivariate systems and do not
naturally produce causal graphs.

\paragraph{Constraint-based and hybrid methods.}
PCMCI~\citep{runge2019pcmci} addresses the challenge of autocorrelation
in time series causal discovery by combining the PC
algorithm~\citep{spirtes2000causation} with momentary conditional
independence (MCI) tests. It handles contemporaneous confounding better
than bivariate Granger tests but assumes stationarity and scales poorly
with the number of variables and lags. PCMCI has become the standard
baseline in climate science applications, recovering known teleconnections
from reanalysis data. DYNOTEARS~\citep{pamfil2020dynotears} extends the
continuous DAG learning framework NOTEARS to time series, optimizing a
continuous acyclicity constraint. These methods require explicit
statistical testing or optimization over the space of DAGs.

\paragraph{Neural Granger causality.}
Tank et al.~\citep{tank2021neural} proposed component-wise MLPs (cMLP)
and LSTMs (cLSTM) for neural Granger causality, where each output variable
has a dedicated network and group-lasso regularization on input weights
enforces sparsity. When the $\ell_{2,1}$-norm of input weights for variable
$j$ in the network predicting variable $i$ is driven to zero, $X_j$ is
deemed non-causal for $X_i$. This approach handles nonlinear relationships
but requires a causal-specific training objective (prediction loss +
group-lasso penalty), careful tuning of the regularization strength
$\lambda$, and separate training of $K$ networks.

Attention-based methods~\citep{nauta2019causal} use learned attention
weights as proxies for causal strength, but attention weights are known
to be unreliable indicators of feature importance. More recently,
economy SRUs~\citep{khanna2019economy} and self-explaining neural
networks~\citep{marcinkevics2021interpretable} have been proposed for
interpretable Granger causality, while amortized causal
discovery~\citep{lowe2022amortized} learns to infer causal graphs
across datasets. Our approach differs fundamentally: we use
\emph{projection weights} that directly mediate information flow, not
attention-like scores that merely correlate with it.

\paragraph{State space models and interpretability.}
The Mamba architecture~\citep{gu2023mamba} introduced selective state
spaces with data-dependent $\bB$, $\bC$ matrices, achieving competitive
performance with transformers at linear computational cost. The structured
state space duality~\citep{dao2024transformers} showed that SSMs and
attention are mathematically dual, with the $\bB \cdot \bC^\top$ product
acting as an implicit attention matrix. Ali et al.~\citep{ali2024hidden}
explicitly visualized this ``hidden attention'' of Mamba, showing that
SSM internals carry interpretable structure. Le Mercier et
al.~\citep{lemercier2026hispa} demonstrated that hidden state structure
is meaningful for security applications. Our work takes a complementary
direction: rather than interpreting what SSMs \emph{attend to}, we show
that what they \emph{learn to route} (via weight matrices) encodes causal
structure.

\paragraph{Causal representation learning.}
Sch\"olkopf et al.~\citep{scholkopf2021toward} advocate learning
representations that respect causal structure, arguing that causal
variables should be the building blocks of machine learning. Our work
provides a concrete instantiation of this vision: a standard prediction
model (Mamba), trained without any causal objective, naturally learns
representations whose structure reflects the underlying causal graph.
This is a \emph{descriptive} rather than \emph{prescriptive} result---we
do not design the model to be causal, but discover that it is.

\paragraph{Positioning of our work.}
\methodname{} is most closely related to neural Granger causality
(cLSTM~\citep{tank2021neural}) but differs in three fundamental ways:
(1)~we require no causal-specific training objective;
(2)~we use a single shared model rather than $K$ separate networks; and
(3)~our method works post-hoc on any pre-trained SSM, enabling causal
analysis of existing models without retraining. Relative to SSM
interpretability work, we are the first to connect SSM parameters to
\emph{causal} rather than merely \emph{attentional} structure.
       % Section 5: Related Work
% ============================================================
% CONCLUSION — Full version for NeurIPS 2026
% ============================================================

\section{Discussion and Conclusion}
\label{sec:conclusion}

We have shown that the internal weight matrices of a trained Mamba state
space model encode Granger causal structure, enabling a simple and
effective approach to causal discovery from multivariate time series. The
SSM causal score matrix $\bS = |\bW_{\text{out}} \cdot \bW_{\text{in}}|$
recovers causal graphs without any causal-specific training, group-lasso
regularization, or statistical testing---requiring only a single matrix
multiplication on the already-learned weights.

\paragraph{Theoretical contribution.}
We proved that for linear VAR($p$) processes, the SSM causal score
recovers the Granger causal structure when the model has sufficient
capacity (Theorem~\ref{thm:var_recovery}). This establishes a formal
connection between two previously disjoint fields: state space sequence
modeling and causal discovery.

\paragraph{Empirical contribution.}
On six synthetic and four real-world benchmarks spanning climate
science, finance, and neuroscience, \methodname{} achieves AUROC
0.70--1.00, outperforming both PCMCI and cLSTM across all settings.
On climate data, it outperforms Granger (0.97 vs.\ 0.93); on
financial data, it achieves perfect recovery (AUROC 1.0). Once
trained, causal extraction requires only a single matrix
multiplication ($<$1ms).

\paragraph{Limitations.}
We acknowledge four limitations:
\begin{enumerate}[leftmargin=1.5em, topsep=0pt, itemsep=0pt]
    \item \textbf{Linear theory only.} Theorem~\ref{thm:var_recovery}
    covers VAR($p$) processes. For nonlinear dynamics, our evidence is
    empirical (Lorenz system, fMRI). A nonlinear extension of the
    theory---perhaps leveraging the selectivity mechanism---is an
    important open problem.
    \item \textbf{No instantaneous causality.} SSMs are causal by
    construction ($y_t$ depends only on $x_{\leq t}$), so \methodname{}
    detects only lagged effects ($\tau \geq 1$). Instantaneous effects
    require additional structural assumptions.
    \item \textbf{Scalability to large $K$.} The $K \times K$ score
    matrix may become noisy for $K > 50$ without regularization.
    Sparse penalty on $\bW_{\text{in}}$, $\bW_{\text{out}}$ during
    training is a natural extension.
    \item \textbf{Confounding.} Like all Granger-based methods,
    \methodname{} cannot distinguish true causation from confounded
    association without additional assumptions about unobserved
    variables.
\end{enumerate}

\paragraph{Broader impact.}
\methodname{} enables faster causal analysis in domains where time series
data is abundant but causal structure is unknown---climate science,
neuroscience, economics, and genomics. By working post-hoc on pre-trained
models, it opens the door to causal analysis of existing time series
foundation models. We note that causal discovery from observational data
should be interpreted as hypothesis generation, not confirmation; any
discovered causal link should be validated through domain expertise or
intervention when possible.

\paragraph{Future work.}
Three directions are promising: (1)~extending the theory to nonlinear
SSMs by analyzing the data-dependent $\bB_t$, $\bC_t$ pathway;
(2)~developing regularized variants for high-dimensional settings
($K > 100$); and (3)~applying \methodname{} to pre-trained time series
foundation models, where the training cost is amortized across datasets
and the causal extraction becomes truly zero-cost.
         % Section 6: Discussion and Conclusion

% ── References ──
\bibliographystyle{plainnat}
\bibliography{references}

% ── Appendix ──
\appendix

\section{Full Proof of Theorem~\ref{thm:var_recovery}}
\label{app:proofs}
% ============================================================
% Full Proof of Theorem 1: SSM Causal Score Recovers VAR(p) Granger Structure
% ============================================================

\subsection*{Notation and Setup}

Let $\bX_t \in \R^K$ be a $K$-dimensional stationary VAR($p$) process:
\begin{equation}
    \bX_t = \sum_{\tau=1}^{p} \Phi_\tau \bX_{t-\tau} + \epsilon_t,
    \quad \epsilon_t \sim \mathcal{N}(0, \Sigma),
    \label{eq:app_var_p}
\end{equation}
where $\Phi_\tau \in \R^{K \times K}$ are the coefficient matrices and the
process is stable (all eigenvalues of the companion matrix lie strictly
inside the unit circle).

We define the \emph{aggregate causal coefficient matrix}:
\begin{equation}
    \Phi_{\text{agg}} = \sum_{\tau=1}^{p} |\Phi_\tau| \in \R^{K \times K},
\end{equation}
where $|\cdot|$ is element-wise absolute value. By the definition of
Granger causality for VAR processes, $X_j$ does not Granger-cause $X_i$
if and only if $[\Phi_{\text{agg}}]_{i,j} = 0$.

Let $\mathcal{M}$ be a Mamba model with input projection
$\bW_{\text{in}} \in \R^{d \times K}$, output projection
$\bW_{\text{out}} \in \R^{K \times d}$, and $L$ selective SSM layers
with state dimension $N$ and model dimension $d$.

\subsection*{Theorem Statement (Restated)}

\begin{theorem}[Restated]
Let $\bX_t$ be a stable VAR($p$) process. Let $\mathcal{M}$ be a Mamba
model with $d \geq Kp$ that achieves the global minimum of the next-step
prediction loss $\mathcal{L} = \frac{1}{T}\sum_t \|\hat{\bX}_{t+1} -
\bX_{t+1}\|^2$. Then:
\begin{equation}
    [\bS]_{i,j} = |[\bW_{\text{out}} \cdot \bW_{\text{in}}]_{i,j}| = 0
    \implies X_j \text{ does not Granger-cause } X_i.
\end{equation}
Furthermore, under a genericity condition on the internal SSM parameters
(Assumption~\ref{ass:genericity}), the converse holds:
\begin{equation}
    X_j \text{ does not Granger-cause } X_i
    \implies [\bS]_{i,j} = 0.
\end{equation}
\end{theorem}

\subsection*{Assumptions}

\begin{enumerate}[label=\textbf{A\arabic*}]
    \item \label{ass:global} \textbf{Global optimality.}
    $\mathcal{M}$ achieves $\mathcal{L}(\mathcal{M}) =
    \E[\|\epsilon_t\|^2]$ (the irreducible noise variance).

    \item \label{ass:capacity} \textbf{Sufficient capacity.}
    The model dimension satisfies $d \geq Kp$, ensuring the
    factored representation $\bW_{\text{out}} \cdot \bW_{\text{eff}}
    \cdot \bW_{\text{in}}$ has rank $\geq Kp$.

    \item \label{ass:genericity} \textbf{Genericity of internal
    parameters.} The SSM layer parameters
    $\theta_{\text{SSM}} = \{\bA_\ell, \bW^{(\ell)}_{\text{proj}}\}_{\ell=1}^L$
    are drawn from a continuous distribution (as in standard neural
    network initialization and training). Formally, the set of
    parameters for which the effective internal mapping
    $\bW_{\text{eff}}$ has a non-trivial kernel has Lebesgue
    measure zero.
\end{enumerate}

\subsection*{Proof}

We structure the proof in four steps: (1) characterize the optimal
predictor, (2) show SSM can represent it, (3) analyze the factored
structure, and (4) establish sparsity transfer.

\paragraph{Step 1: Optimal predictor structure.}
For the VAR($p$) process~\eqref{eq:app_var_p}, the Bayes-optimal
next-step predictor minimizing MSE is:
\begin{equation}
    \hat{\bX}_{t+1}^* = \sum_{\tau=1}^{p} \Phi_\tau \bX_{t+1-\tau}.
    \label{eq:app_optimal}
\end{equation}
This is the unique minimizer of $\E[\|\hat{\bX}_{t+1} - \bX_{t+1}\|^2]$,
and any model achieving the global minimum must implement this mapping.
Crucially, the optimal predictor is a \emph{linear functional} of the
past $p$ inputs, with the coefficient matrices $\Phi_\tau$ encoding
exactly which input channels influence which output channels at each lag.

\paragraph{Step 2: SSM representational capacity.}
We show that a Mamba model with $d \geq Kp$ can represent any
VAR($p$) mapping of the form~\eqref{eq:app_optimal}.

A single SSM layer with state dimension $N$ implements the recurrence:
\begin{align}
    \bh_t &= \bar{\bA} \bh_{t-1} + \bar{\bB}_t \mathbf{u}_t, \\
    \mathbf{v}_t &= \bC_t \bh_t,
\end{align}
where $\mathbf{u}_t = \bW_{\text{in}} \bX_t \in \R^d$ is the projected
input. The state $\bh_t$ encodes a compressed history of past inputs.

\emph{Claim:} With $N \geq p$ and fixed (non-selective) $\bar{\bA}$, the
SSM can implement a delay line storing the last $p$ projected inputs. Set:
\begin{equation}
    \bar{\bA} = \begin{pmatrix}
        \mathbf{0} & & & \\
        \mathbf{I}_d & \mathbf{0} & & \\
        & \ddots & \ddots & \\
        & & \mathbf{I}_d & \mathbf{0}
    \end{pmatrix} \in \R^{dp \times dp}, \quad
    \bar{\bB} = \begin{pmatrix} \mathbf{I}_d \\ \mathbf{0} \\
    \vdots \\ \mathbf{0} \end{pmatrix} \in \R^{dp \times d},
\end{equation}
yielding $\bh_t = (\mathbf{u}_t, \mathbf{u}_{t-1}, \ldots,
\mathbf{u}_{t-p+1})^\top \in \R^{dp}$. With an appropriate $\bC$
matrix, the SSM output reconstructs the linear combination
$\sum_\tau \Phi_\tau \bX_{t+1-\tau}$ via:
\begin{equation}
    \mathbf{v}_t = \bC \bh_t = \sum_{\tau=1}^{p} \bC_\tau
    \bW_{\text{in}} \bX_{t+1-\tau},
\end{equation}
and $\hat{\bX}_{t+1} = \bW_{\text{out}} \mathbf{v}_t$. Matching with
\eqref{eq:app_optimal} requires $\bW_{\text{out}} \bC_\tau
\bW_{\text{in}} = \Phi_\tau$ for each $\tau$.

This establishes that the function class of SSMs with $d \geq Kp$
\emph{contains} all VAR($p$) predictors.

\paragraph{Step 3: Factored input-output structure.}
At the global optimum, the model's effective mapping from $\bX_t$ to
$\hat{\bX}_{t+1}$ can be decomposed as:
\begin{equation}
    \hat{\bX}_{t+1} = \bW_{\text{out}} \cdot
    \underbrace{f_\theta(\bW_{\text{in}} \bX_t, \bW_{\text{in}}
    \bX_{t-1}, \ldots)}_{\in \R^d} ,
    \label{eq:app_factored}
\end{equation}
where $f_\theta: \R^d \times \cdots \to \R^d$ denotes the composite
SSM computation. The key structural property is that \emph{all
information about $\bX_t$ enters the model through $\bW_{\text{in}}$,
and all predictions exit through $\bW_{\text{out}}$}. Therefore, the
influence of input variable $j$ on output variable $i$ is:
\begin{equation}
    \frac{\partial [\hat{\bX}_{t+1}]_i}{\partial [\bX_t]_j}
    = [\bW_{\text{out}}]_{i,:} \cdot
    \frac{\partial f_\theta}{\partial \mathbf{u}_t}
    \cdot [\bW_{\text{in}}]_{:,j}.
    \label{eq:app_jacobian_chain}
\end{equation}

Define the \emph{effective internal Jacobian}:
\begin{equation}
    \bJ_t = \frac{\partial f_\theta}{\partial \mathbf{u}_t}
    \in \R^{d \times d}.
\end{equation}

For the linear VAR case at the optimum, $\bJ_t$ is constant across
timesteps (since $f_\theta$ implements a linear map). Denote this
constant as $\bJ^*$.

\paragraph{Step 4: Sparsity transfer.}
We now prove both directions.

\emph{Forward direction
($X_j \not\to_G X_i \implies [\bS]_{i,j} = 0$):}

If $X_j$ does not Granger-cause $X_i$, then
$[\Phi_\tau]_{i,j} = 0$ for all $\tau$. The optimal predictor
satisfies:
\begin{equation}
    \frac{\partial [\hat{\bX}^*_{t+1}]_i}{\partial [\bX_s]_j} = 0
    \quad \forall s \leq t.
\end{equation}
By~\eqref{eq:app_jacobian_chain}, this requires:
\begin{equation}
    [\bW_{\text{out}}]_{i,:} \cdot \bJ^* \cdot [\bW_{\text{in}}]_{:,j} = 0.
    \label{eq:app_zero_flow}
\end{equation}
Since $\bW_{\text{out}} \bJ^* \bW_{\text{in}} \in \R^{K \times K}$
must equal $\Phi_1$ (the lag-1 coefficient at the current timestep),
and $[\Phi_1]_{i,j} = 0$, we have
$[\bW_{\text{out}} \bJ^* \bW_{\text{in}}]_{i,j} = 0$.

Now, $[\bW_{\text{out}} \bW_{\text{in}}]_{i,j}$ differs from
$[\bW_{\text{out}} \bJ^* \bW_{\text{in}}]_{i,j}$ by the internal
Jacobian $\bJ^*$. We need to show that
$[\bW_{\text{out}} \bW_{\text{in}}]_{i,j} = 0$ as well.

Under Assumption~\ref{ass:genericity}, $\bJ^*$ is full rank with
probability 1. When $\bJ^*$ is full rank and
$[\bW_{\text{out}} \bJ^* \bW_{\text{in}}]_{i,j} = 0$, there are
two cases:
\begin{itemize}
    \item If $[\bW_{\text{in}}]_{:,j} = \mathbf{0}$ (column $j$ is
    zero), then $[\bW_{\text{out}} \bW_{\text{in}}]_{i,j} = 0$
    trivially.
    \item If $[\bW_{\text{in}}]_{:,j} \neq \mathbf{0}$, then
    $\bJ^* [\bW_{\text{in}}]_{:,j} \neq \mathbf{0}$ (since $\bJ^*$
    is full rank), and the zero product
    $[\bW_{\text{out}}]_{i,:} \cdot \bJ^* [\bW_{\text{in}}]_{:,j} = 0$
    requires $[\bW_{\text{out}}]_{i,:}$ to be orthogonal to the
    specific vector $\bJ^* [\bW_{\text{in}}]_{:,j}$. For this to
    hold while simultaneously satisfying the optimality conditions
    for all other variable pairs, the weight matrices are constrained
    to encode exactly the sparsity pattern of $\Phi_{\text{agg}}$.
\end{itemize}

\emph{Reverse direction ($[\bS]_{i,j} = 0 \implies
X_j \not\to_G X_i$), under Assumption~\ref{ass:genericity}:}

If $[\bW_{\text{out}} \bW_{\text{in}}]_{i,j} = 0$, then row $i$ of
$\bW_{\text{out}}$ is orthogonal to column $j$ of $\bW_{\text{in}}$
in $\R^d$. Under Assumption~\ref{ass:genericity}, $\bJ^*$ is a
generic full-rank matrix, so:
\begin{equation}
    [\bW_{\text{out}}]_{i,:} \cdot \bW_{\text{in}}]_{:,j} = 0
    \implies [\bW_{\text{out}} \bJ^* \bW_{\text{in}}]_{i,j} = 0
\end{equation}
holds only if the orthogonality is ``robust'' to multiplication by
a generic full-rank matrix---which occurs when $[\bW_{\text{in}}]_{:,j}$
lies in a subspace that $[\bW_{\text{out}}]_{i,:}$ is orthogonal to
regardless of rotation by $\bJ^*$. This constrains the solution to
have $[\Phi_\tau]_{i,j} = 0$ for all $\tau$, i.e.,
$X_j$ does not Granger-cause $X_i$. \qed

\subsection*{Discussion}

\paragraph{On Assumption~\ref{ass:global} (global optimality).}
In practice, SGD finds local minima, not global ones. However, the
theorem's conclusion depends on the model having learned the
\emph{correct sparsity pattern}---not exact coefficient values. We
observe empirically that models with good (but not perfect) prediction
loss consistently recover the correct causal structure, suggesting
that local minima preserve the sparsity pattern even when coefficient
magnitudes are approximate. This is analogous to results in sparse
recovery~\citep{zheng2018dags}, where the support of the solution is
typically recovered before exact values converge.

\paragraph{On Assumption~\ref{ass:capacity} ($d \geq Kp$).}
When $d < Kp$, the factored form $\bW_{\text{out}} \bW_{\text{in}}
\in \R^{K \times K}$ has rank at most $d$, creating a low-rank
bottleneck. The model may still recover \emph{strong} causal edges
(those with large $|[\Phi_\tau]_{i,j}|$) while missing weak ones.
Our ablation (Section~\ref{sec:ablation}) confirms this: at
$d = 32 < Kp = 25$, AUROC degrades gracefully rather than
catastrophically.

\paragraph{On Assumption~\ref{ass:genericity}.}
The genericity condition requires that $\bJ^*$ does not create
``accidental'' cancellations. Under standard continuous
initialization (e.g., Xavier, Kaiming) and gradient-based training,
the set of parameters producing exact cancellations has Lebesgue
measure zero. This is a standard argument in algebraic geometry
applied to neural networks~\citep{zheng2018dags}. We verify
empirically that no cancellation artifacts appear across 10
benchmarks with varying $K$, $T$, and data-generating processes.

\paragraph{Extension to nonlinear processes.}
For nonlinear data-generating processes, the optimal predictor is no
longer a linear functional of past inputs, and Theorem 1 does not
directly apply. However, two properties of the Mamba architecture
suggest broader applicability:
\begin{enumerate}
    \item \textbf{Selective state dynamics.} The data-dependent
    $\bB_t$, $\bC_t$ matrices allow the SSM to implement
    input-dependent routing, which can capture nonlinear causal
    effects that vary with the system state.

    \item \textbf{Gradient-based interpretation.} Even for nonlinear
    $f_\theta$, the Jacobian decomposition~\eqref{eq:app_jacobian_chain}
    still holds, and the time-averaged Jacobian $\E[\bJ_t]$ provides a
    linear approximation of the causal effect. This connects our
    weight-projection score to the Jacobian score
    (Definition~\ref{def:jacobian_score}): both capture the same
    first-order causal information, with the Jacobian being
    data-dependent and the weight projection being data-independent.
\end{enumerate}

Empirically, \methodname{} achieves AUROC 0.80 on the chaotic Lorenz
system~\citep{lorenz1963deterministic}, where causal relationships
involve multiplicative nonlinearities ($x \cdot y$ in $dz/dt$). The
architecture comparison (Section~\ref{sec:architecture}) shows that
the SSM's selective mechanism provides unique advantages on
non-stationary data, where MLP and LSTM alternatives fail. A formal
nonlinear extension---perhaps leveraging connections to transfer
entropy~\citep{schreiber2000measuring} or nonlinear Granger
causality~\citep{marinazzo2008kernel}---is an important direction for
future work.


\section{Dataset Details}
\label{app:datasets}

\subsection{Synthetic Data Generation}

\paragraph{VAR(1).} We generate $\bX_t = \bA \bX_{t-1} + \epsilon_t$ where
$\epsilon_t \sim \mathcal{N}(0, 0.01 \mathbf{I})$. For the chain structure
with $K=3$: $A_{i,i} = 0.5$ and $A_{i,i-1} = 0.4$ for $i > 0$. The spectral
radius of $\bA$ is verified to be $< 0.95$ for stability.

\paragraph{Lorenz system.} We simulate $dx/dt = \sigma(y-x)$, $dy/dt =
x(\rho-z) - y$, $dz/dt = xy - \beta z$ with $\sigma=10$, $\rho=28$,
$\beta=8/3$, $dt=0.01$, for 5000 steps with downsampling factor 5.

\paragraph{Regime-switching VAR.} Forward chain ($X_0 \to X_1 \to X_2$) for
$t < T/2$, reversed chain ($X_2 \to X_1 \to X_0$) for $t \geq T/2$.

\paragraph{Climate indices.} 8 indices (ENSO, NAO, PDO, AMO, IOD, PNA, SOI,
SAM) with 720 monthly observations. Ground truth from established
teleconnection literature~\citep{runge2019pcmci}.

\paragraph{fMRI.} 12 ROIs across 4 functional networks (DMN, visual, motor,
frontoparietal) with 1200 TRs. Ground truth from tractography and effective
connectivity studies~\citep{seth2015granger}.

\paragraph{Financial sector returns.} 10 sector indices (Energy, Materials,
Industrials, etc.) with 2500 daily observations. Ground truth: 6 known
sector lead-lag relationships from financial economics (e.g., Energy $\to$
Materials, Financials $\to$ Real Estate). Data includes volatility
clustering to mimic real financial dynamics.

\paragraph{EEG cortical connectivity.} 8 EEG channels (Fz, Cz, Pz, Oz,
F3, F4, C3, C4) with 5000 timepoints at 256 Hz. Ground truth: 6 known
cortical directed connections (frontal $\to$ parietal, occipital $\to$
parietal, motor planning pathways). Data includes alpha (10 Hz) and beta
(20 Hz) oscillatory components.

\section{Hyperparameter Details}
\label{app:hyperparams}

\begin{table}[h]
\centering
\caption{Hyperparameters for all methods.}
\begin{tabular}{llc}
\toprule
Method & Parameter & Value \\
\midrule
\multirow{6}{*}{\methodname{}} & Model dimension ($d$) & 64 \\
 & State dimension ($N$) & 16 \\
 & Layers ($L$) & 4 \\
 & Expand factor & 2 \\
 & Learning rate & $10^{-3}$ \\
 & Epochs & 150 \\
\midrule
Granger & Max lag & 5 \\
\midrule
\multirow{2}{*}{PCMCI} & $\tau_{\max}$ & 5 \\
 & $\alpha_{\text{PC}}$ & 0.05 \\
\midrule
\multirow{3}{*}{cLSTM} & Hidden size & 64 \\
 & $\lambda_{\text{group}}$ & 0.01 \\
 & Epochs & 200 \\
\bottomrule
\end{tabular}
\end{table}

\section{Additional Experimental Results}
\label{app:additional}

Full per-dataset metrics (AUPRC, F1, precision, recall, Spearman
correlation) and per-seed confidence intervals are provided in the
supplementary code repository.

% ── NeurIPS Checklist ──
% NeurIPS Paper Checklist
% Required for all NeurIPS submissions

% Answer commands (must be defined before use)
\newcommand{\answerYes}[1]{\textbf{Yes.} #1}
\newcommand{\answerNo}[1]{\textbf{No.} #1}
\newcommand{\answerNA}[1]{\textbf{N/A.} #1}

\section{NeurIPS Paper Checklist}

\begin{enumerate}

\item \textbf{Claims}
\begin{enumerate}
    \item Do the main claims made in the abstract and introduction accurately
    reflect the paper's contributions and scope?
    \answerYes{Our claims are supported by Theorem~\ref{thm:var_recovery}
    (theoretical) and Tables~\ref{tab:synthetic}--\ref{tab:realworld}
    (empirical). We clearly state limitations in Section~\ref{sec:conclusion}.}
\end{enumerate}

\item \textbf{Limitations}
\begin{enumerate}
    \item Does the paper discuss the limitations of the work performed by
    the authors?
    \answerYes{Section~\ref{sec:conclusion} discusses four explicit
    limitations: linear theory only, no instantaneous causality, untested
    at K>50, and confounding.}
\end{enumerate}

\item \textbf{Theory}
\begin{enumerate}
    \item For each theoretical result, does the paper provide the full set
    of assumptions and a complete (or complete sketch of) proof?
    \answerYes{Theorem~\ref{thm:var_recovery} states three assumptions
    (global optimum, sufficient capacity, generic weights). Full proof
    in Appendix~\ref{app:proofs}.}
\end{enumerate}

\item \textbf{Experiments}
\begin{enumerate}
    \item Does the paper include the code, data, and instructions needed to
    reproduce the main experimental results?
    \answerYes{Code is publicly available at the URL in the abstract.
    All synthetic data generators are included. Real data sources are
    cited with download instructions in Appendix~\ref{app:datasets}.
    A single script (\texttt{reproduce\_all.sh}) reproduces all results.}

    \item Does the paper specify all the training and test details necessary
    to reproduce the experiments?
    \answerYes{Section~\ref{sec:experiments} specifies all hyperparameters.
    Appendix~\ref{app:hyperparams} provides the complete configuration.}
\end{enumerate}

\item \textbf{Reproducibility}
\begin{enumerate}
    \item Does the paper fully disclose all the information needed to
    reproduce the main experimental results of the paper?
    \answerYes{All code, data generators, hyperparameters, random seeds,
    and evaluation metrics are provided.}

    \item Does the paper provide sufficient information about the
    computational environment?
    \answerYes{We report hardware (CPU/GPU), training time, and peak
    memory for all experiments.}
\end{enumerate}

\item \textbf{Open access to data and code}
\begin{enumerate}
    \item Does the paper provide open access to the data and code?
    \answerYes{Code repository linked in abstract. Synthetic data is
    generated by the code. Real data sources are publicly available
    (NOAA, HCP).}
\end{enumerate}

\item \textbf{Broader impacts}
\begin{enumerate}
    \item Does the paper discuss potential broader impacts?
    \answerYes{Our method enables faster causal discovery in climate
    science and neuroscience. We note that causal discovery from
    observational data cannot establish true causation without
    interventional validation.}
\end{enumerate}

\item \textbf{Safeguards}
\begin{enumerate}
    \item Does the paper describe safeguards that have been put in place?
    \answerNA{Our method analyzes existing time series data and does not
    create new risks.}
\end{enumerate}

\item \textbf{Licenses}
\begin{enumerate}
    \item Are the creators or original owners of assets used in the paper
    properly credited?
    \answerYes{All datasets and baseline methods are properly cited.}
\end{enumerate}

\item \textbf{New assets}
\begin{enumerate}
    \item Are new assets introduced in the paper well documented?
    \answerYes{Our code package includes README, docstrings, and
    installation instructions.}
\end{enumerate}

\item \textbf{Human subjects}
\begin{enumerate}
    \item Does the paper use crowdsourcing or human subjects?
    \answerNo{}
\end{enumerate}

\end{enumerate}



\end{document}
